% GENERATED FILE. Edit canonical lessons, then run npm run generate:books.
% canonical-source-hash: fnv1a64:eeb355393a367524
% canonical-lessons: MR-C179-nyayalay, MR-C179-nyayadhis, MR-C179-mahapaur, MR-C179-nivadnuk, MR-C179-rajkaran, MR-R179-first-pass-words-from-chapters-171-174, MR-R179-second-pass-words-from-chapters-175-179

\chapter{Court, Judge, Mayor, Election, Politics}
\label{ch:mr-court-judge-mayor-election-politics}
\hlchaptermodality{179}

\begin{chapteropening}
\textbf{By the end of this chapter:} \emph{I can say a court, a judge, a mayor, an election and politics.}

Complete the last lesson of chapter 179: I can say a court, a judge, a mayor, an election and politics.
\end{chapteropening}

\section[\texorpdfstring{\mr{न्यायालय} (nyāyālay)}{न्यायालय (nyāyālay)}]{\mr{न्यायालय} (nyāyālay) — a court}

\label{lesson:MR-C179-nyayalay}



\begin{quote}
Before the new one: say the Marathi for a bell, then the Marathi for a compartment.
\end{quote}

\subsection*{You'll want to know: \mr{न्यायालय}}
\textbf{\mr{न्यायालय}} — \emph{nyāyālay} — \textquotedblleft{}a court\textquotedblright{}.

\begin{grammarlens}[title={What you've built}]
One more word for this chapter.
\end{grammarlens}

\subsection*{Your turn}
\begin{itemize}
\raggedright
  \item \emph{Say it:} \emph{nyāyālay}
  \item \emph{Say it:} \emph{nyāyālay}, once more
  \item \emph{Say it:} say \emph{nyāyālay}
  \item \emph{Recall:} say the Marathi for a bell, then the Marathi for a compartment, then say \emph{nyāyālay} again
\end{itemize}

\begin{tcolorbox}[breakable,colback=teal!4,colframe=teal!35!black,title={Before you move on}]
What does \mr{न्यायालय} mean? (\textquotedblleft{}A court\textquotedblright{}.) Say it once more.
\end{tcolorbox}
\section[\texorpdfstring{\mr{न्यायाधीश} (nyāyādhīś)}{न्यायाधीश (nyāyādhīś)}]{\mr{न्यायाधीश} (nyāyādhīś) — a judge}

\label{lesson:MR-C179-nyayadhis}



\begin{quote}
Before the new one: say the Marathi for a compartment, then the Marathi for a court.
\end{quote}

\subsection*{You'll want to know: \mr{न्यायाधीश}}
\textbf{\mr{न्यायाधीश}} — \emph{nyāyādhīś} — \textquotedblleft{}a judge\textquotedblright{}.

\begin{grammarlens}[title={What you've built}]
One more word for this chapter.
\end{grammarlens}

\subsection*{Your turn}
\begin{itemize}
\raggedright
  \item \emph{Say it:} \emph{nyāyādhīś}
  \item \emph{Say it:} \emph{nyāyādhīś}, once more
  \item \emph{Say it:} say \emph{nyāyādhīś}
  \item \emph{Recall:} say the Marathi for a compartment, then the Marathi for a court, then say \emph{nyāyādhīś} again
\end{itemize}

\begin{tcolorbox}[breakable,colback=teal!4,colframe=teal!35!black,title={Before you move on}]
What does \mr{न्यायाधीश} mean? (\textquotedblleft{}A judge\textquotedblright{}.) Say it once more.
\end{tcolorbox}
\section[\texorpdfstring{\mr{महापौर} (mahāpaur)}{महापौर (mahāpaur)}]{\mr{महापौर} (mahāpaur) — a mayor}

\label{lesson:MR-C179-mahapaur}



\begin{quote}
Before the new one: say the Marathi for a court, then the Marathi for a judge.
\end{quote}

\subsection*{You'll want to know: \mr{महापौर}}
\textbf{\mr{महापौर}} — \emph{mahāpaur} — \textquotedblleft{}a mayor\textquotedblright{}.

\begin{grammarlens}[title={What you've built}]
One more word for this chapter.
\end{grammarlens}

\subsection*{Your turn}
\begin{itemize}
\raggedright
  \item \emph{Say it:} \emph{mahāpaur}
  \item \emph{Say it:} \emph{mahāpaur}, once more
  \item \emph{Say it:} say \emph{mahāpaur}
  \item \emph{Recall:} say the Marathi for a court, then the Marathi for a judge, then say \emph{mahāpaur} again
\end{itemize}

\begin{tcolorbox}[breakable,colback=teal!4,colframe=teal!35!black,title={Before you move on}]
What does \mr{महापौर} mean? (\textquotedblleft{}A mayor\textquotedblright{}.) Say it once more.
\end{tcolorbox}
\section[\texorpdfstring{\mr{निवडणूक} (nivaḍṇūk)}{निवडणूक (nivaḍṇūk)}]{\mr{निवडणूक} (nivaḍṇūk) — an election}

\label{lesson:MR-C179-nivadnuk}



\begin{quote}
Before the new one: say the Marathi for a judge, then the Marathi for a mayor.
\end{quote}

\subsection*{You'll want to know: \mr{निवडणूक}}
\textbf{\mr{निवडणूक}} — \emph{nivaḍṇūk} — \textquotedblleft{}an election\textquotedblright{}.

\begin{grammarlens}[title={What you've built}]
One more word for this chapter.
\end{grammarlens}

\subsection*{Your turn}
\begin{itemize}
\raggedright
  \item \emph{Say it:} \emph{nivaḍṇūk}
  \item \emph{Say it:} \emph{nivaḍṇūk}, once more
  \item \emph{Say it:} say \emph{nivaḍṇūk}
  \item \emph{Recall:} say the Marathi for a judge, then the Marathi for a mayor, then say \emph{nivaḍṇūk} again
\end{itemize}

\begin{tcolorbox}[breakable,colback=teal!4,colframe=teal!35!black,title={Before you move on}]
What does \mr{निवडणूक} mean? (\textquotedblleft{}An election\textquotedblright{}.) Say it once more.
\end{tcolorbox}
\section[\texorpdfstring{\mr{राजकारण} (rājkāraṇ)}{राजकारण (rājkāraṇ)}]{\mr{राजकारण} (rājkāraṇ) — politics}

\label{lesson:MR-C179-rajkaran}



\begin{quote}
Before the new one: say the Marathi for a mayor, then the Marathi for an election.
\end{quote}

\subsection*{You'll want to know: \mr{राजकारण}}
\textbf{\mr{राजकारण}} — \emph{rājkāraṇ} — \textquotedblleft{}politics\textquotedblright{}.

\begin{grammarlens}[title={What you've built}]
One more word for this chapter.
\end{grammarlens}

\subsection*{Your turn}
\begin{itemize}
\raggedright
  \item \emph{Say it:} \emph{rājkāraṇ}
  \item \emph{Say it:} \emph{rājkāraṇ}, once more
  \item \emph{Say it:} say \emph{rājkāraṇ}
  \item \emph{Recall:} say the Marathi for a mayor, then the Marathi for an election, then say \emph{rājkāraṇ} again
\end{itemize}

\begin{tcolorbox}[breakable,colback=teal!4,colframe=teal!35!black,title={Before you move on}]
What does \mr{राजकारण} mean? (\textquotedblleft{}Politics\textquotedblright{}.) Say it once more.
\end{tcolorbox}
\section[(dialogue)]{Words from chapters 171-174 — a first pass}

\label{lesson:MR-R179-first-pass-words-from-chapters-171-174}



\begin{quote}
Say the words for an election and politics.
\end{quote}

\subsection*{Your turn}
Say these aloud:

\begin{itemize}
\raggedright
  \item rarely, then, finally, still, suddenly
  \item around, everywhere, elsewhere, over there, in the middle
  \item even so, although, despite, unfortunately, luckily
  \item comfortable, polite, rude, unpleasant, enthusiastic
\end{itemize}

\begin{tcolorbox}[breakable,colback=teal!4,colframe=teal!35!black,title={Before you move on}]
Rarely? (\textbf{\mr{क्वचित}}, \emph{kvacit}.) Then? (\textbf{\mr{मग}}, \emph{mag}.) Finally? (\textbf{\mr{शेवटी}}, \emph{śevṭī}.) Still? (\textbf{\mr{अजून}}, \emph{ajūn}.) Suddenly? (\textbf{\mr{अचानक}}, \emph{acānak}.) Around? (\textbf{\mr{भोवती}}, \emph{bhovtī}.) Everywhere? (\textbf{\mr{सगळीकडे}}, \emph{sagḷīkaḍe}.) Elsewhere? (\textbf{\mr{इतरत्र}}, \emph{itaratra}.) Over there? (\textbf{\mr{तिकडे}}, \emph{tikaḍe}.) In the middle? (\textbf{\mr{मध्यभागी}}, \emph{madhyabhāgī}.) Even so? (\textbf{\mr{तरीही}}, \emph{tarīhī}.) Although? (\textbf{\mr{जरी}}, \emph{jarī}.) Despite? (\textbf{\mr{असूनही}}, \emph{asūnhī}.) Unfortunately? (\textbf{\mr{दुर्दैवाने}}, \emph{durdaivāne}.) Luckily? (\textbf{\mr{सुदैवाने}}, \emph{sudaivāne}.) Comfortable? (\textbf{\mr{आरामदायक}}, \emph{ārāmdāyak}.) Polite? (\textbf{\mr{सभ्य}}, \emph{sabhya}.) Rude? (\textbf{\mr{उद्धट}}, \emph{uddhaṭ}.) Unpleasant? (\textbf{\mr{अप्रिय}}, \emph{apriya}.) Enthusiastic? (\textbf{\mr{उत्साही}}, \emph{utsāhī}.)
\end{tcolorbox}
\section[(dialogue)]{Words from chapters 175-179 — a second pass}

\label{lesson:MR-R179-second-pass-words-from-chapters-175-179}



\begin{quote}
Say the words for strange and different.
\end{quote}

\subsection*{Your turn}
Say these aloud:

\begin{itemize}
\raggedright
  \item strange, different, similar, proper, wrong
  \item cloudy, rainy, possible, impossible, main
  \item affectionate, kind, cruel, humble, arrogant
  \item a flat, a dining room, a balcony, a bell, a compartment
  \item a court, a judge, a mayor, an election, politics
\end{itemize}

\begin{tcolorbox}[breakable,colback=teal!4,colframe=teal!35!black,title={Before you move on}]
Strange? (\textbf{\mr{विचित्र}}, \emph{vicitra}.) Different? (\textbf{\mr{वेगळा} / \mr{वेगळी}}, \emph{vegḷā / vegḷī}.) Similar? (\textbf{\mr{सारखा} / \mr{सारखी}}, \emph{sārkhā / sārkhī}.) Proper? (\textbf{\mr{योग्य}}, \emph{yogya}.) Wrong? (\textbf{\mr{अयोग्य}}, \emph{ayogya}.) Cloudy? (\textbf{\mr{ढगाळ}}, \emph{ḍhagāḷ}.) Rainy? (\textbf{\mr{पावसाळी}}, \emph{pāvsāḷī}.) Possible? (\textbf{\mr{शक्य}}, \emph{śakya}.) Impossible? (\textbf{\mr{अशक्य}}, \emph{aśakya}.) Main? (\textbf{\mr{मुख्य}}, \emph{mukhya}.) Affectionate? (\textbf{\mr{प्रेमळ}}, \emph{premaḷ}.) Kind? (\textbf{\mr{दयाळू}}, \emph{dayāḷū}.) Cruel? (\textbf{\mr{क्रूर}}, \emph{krūr}.) Humble? (\textbf{\mr{नम्र}}, \emph{namra}.) Arrogant? (\textbf{\mr{गर्विष्ठ}}, \emph{garviṣṭha}.) A flat? (\textbf{\mr{सदनिका}}, \emph{sadnikā}.) A dining room? (\textbf{\mr{जेवणघर}}, \emph{jevaṇghar}.) A balcony? (\textbf{\mr{सज्जा}}, \emph{sajjā}.) A bell? (\textbf{\mr{घंटा}}, \emph{ghaṇṭā}.) A compartment? (\textbf{\mr{कप्पा}}, \emph{kappā}.) A court? (\textbf{\mr{न्यायालय}}, \emph{nyāyālay}.) A judge? (\textbf{\mr{न्यायाधीश}}, \emph{nyāyādhīś}.) A mayor? (\textbf{\mr{महापौर}}, \emph{mahāpaur}.) An election? (\textbf{\mr{निवडणूक}}, \emph{nivaḍṇūk}.) Politics? (\textbf{\mr{राजकारण}}, \emph{rājkāraṇ}.)
\end{tcolorbox}
